\documentclass{ou}
\begin{document}

\thispagestyle{empty}
\begin{center}
\shortstack{\FGAVO}\vspace{4\baselineskip}
\textbf{РЕФЕРАТ}\\
\setstretch{2}
Информационные системы и~технологии\\
\setstretch{1}
Основы архитектуры информационных систем и администрирования\vspace{2\baselineskip}\end{center}
\begin{flushleft}
\doublespacing
\begin{tabularx}{\textwidth}{@{} p{0.4\textwidth} l l @{}}
Руководитель & ст. преподаватель & И. Р. Нигматуллин\\
Студенты     &                   &                  \\
\multirow{-2.27}{\hsize}{\singlespacing Основные компоненты архитектуры информационной системы и их взаимодействие}
             & ГФ25-01Б, 152540046 & А. Д. Балцевич\\
             & ГФ25-01Б, 152537659 & В. Е. Зайцева\\
             & ГФ25-01Б, 152536592 & А. Ю. Бирюкова\\   
\multirow{-2.27}{\hsize}{\singlespacing Роли серверов, клиентов и~сетевых устройств в~архитектуре системы}
             & ГФ25-01Б, 152537684 & О. И. Бочанова\\
             & ГФ25-01Б, 152539986 & М. К. Гинтер\\
             & ГФ25-01Б, 152541471 & А. О. Безруких\\
             & ГФ25-01Б, 152536796 & А. В. Корюкова\\       
\multirow{-2.27}{\hsize}{\singlespacing Функции пользователей, администраторов и операторов в работе информационной системы}
             & ГФ25-01Б, 152537671 & А. Е. Корикова\\
             & ГФ25-01Б, 152536528 & Д. А. Васильева\\
             & ГФ25-01Б, 152540016 & И. Н. Морозова\\  
\end{tabularx}
\end{flushleft}
\vspace*{\fill}
\begin{center}
\city
\end{center}
\newpage
\thispagestyle{empty}
\setstretch{1}
Продолжение титульного~листа реферата по~теме «Основы архитектуры информационных систем и администрирования»
\begin{flushleft}
\doublespacing
\begin{tabularx}{\textwidth}{@{} p{0.4\textwidth} l l @{}}
Студенты     &                     &         \\
\multirow{-2.27}{\hsize}{\singlespacing Принципы организации доступа к информационным ресурсам}
             & ГФ25-01Б, 152541294 & Э. Р. Махмудов\\
             & ГФ25-01Б, 152539829 & В. Д. Бобкова\\
\multirow{-2.27}{\hsize}{\singlespacing Основные задачи системного администрирования и~поддержания работоспособности ИС}
             & ГФ25-01Б, 152536115 & С. В. Идрисова\\
             & ГФ25-01Б, 152540131 & Е. Кичко\\
             & ГФ25-01Б, 152539526 & К. Н. Аркадьева\\
             & ГФ25-01Б, 152536056 & В. А. Карасова\\
\multirow{-2.27}{\hsize}{\singlespacing Инструменты и практики администрирования в~современных вычислительных средах}
             & ГФ25-01Б, 152540765 & А. В. Гайворонская\\
             & ГФ25-01Б, 152540753 & М. И. Бобылев\\
             & ГФ25-01Б, 152539850 & Е. А. Аксаментова\\
\end{tabularx}
\end{flushleft}
\newpage
{
\sloppy
\tableofcontents
}
\newpage
\onehalfspacing
\section{Основные компоненты архитектуры информационной системы и их взаимодействие}
В условиях цифровизации и активного развития информационных технологий информационные системы (ИС) играют ключевую роль в~деятельности организаций, государственных структур и образовательных учреждений. Эффективность и надёжность информационной системы во~многом зависят от её архитектуры. Архитектура информационной системы определяет структуру системы, её основные компоненты и способы их взаимодействия.
\subsection{Понятие архитектуры информационной системы}
Архитектура ИС -- набор значимых решений по поводу организации системы программного обеспечения, набор структурных элементов и их интерфейсов, при помощи которых компонуется система вместе с их поведением, определяемым во взаимодействии между этими элементами, компоновка элементов в постепенно укрупняющиеся подсистемы, а также стиль архитектуры, который направляет эту организацию (элементы и их интерфейсы, взаимодействие и компоновка)\cite{kazan2023}.
\subsection{Основные компоненты архитектуры информационной системы}
Компоненты ИС: отдельные части системы, выполняющие определенные функции (например, базы данных, интерфейсы пользователя)\cite{academia2023}.

Аппаратные компоненты включают в себя технические средства, на~которых функционирует информационная система:
\begin{itemize}[label=-, leftmargin=0em, itemindent=3.35em, nosep]
\item серверы;
\item рабочие станции пользователей;
\item сетевое оборудование;
\item устройства хранения данных.
\end{itemize}

Данный компонент обеспечивает физическую основу работы ИС и~напрямую влияет на её производительность и надёжность.

Программные компоненты представляют совокупность программных средств, обеспечивающих функционирование системы. Они включают:
\begin{itemize}[label=-, leftmargin=0em, itemindent=3.35em, nosep]
\item системное программное обеспечение (операционные системы, СУБД);
\item прикладное программное обеспечение;
\item сервисные и вспомогательные программы.
\end{itemize}

Программное обеспечение реализует бизнес логику системы.

Информационные ресурсы (данные) отвечают за хранение и управление данными, используемыми системой. Здесь могут использоваться базы данных или другие механизмы хранения данных. Уровень хранения данных обеспечивает сохранность информации и обеспечивает доступ к~ней со~стороны других уровней системы\cite{ifmo2024}.

Пользовательский интерфейс отвечает за взаимодействие с пользователем и предоставляет пользовательский интерфейс для ввода и вывода данных. Он может включать в себя вебстраницы, мобильные приложения, графические интерфейсы и другие средства коммуникации с пользователем.

Качество интерфейса напрямую влияет на удобство использования системы и эффективность работы пользователей\cite{ifmo2024}.

Организационное и правовое обеспечение включает регламенты, инструкции, правила эксплуатации и нормативно правовые документы. Он определяет порядок работы с системой и ответственность пользователей.
\subsection{Взаимодействие компонентов архитектуры ИС}
Взаимодействие компонентов архитектуры информационной системы осуществляется по принципу разделения функций. Пользовательский интерфейс передаёт запросы к уровню бизнес логики, который обрабатывает их и обращается к базе данных. Аппаратное и программное обеспечение обеспечивают техническую возможность выполнения этих процессов.

В распределенной архитектуре каждый компонент системы может располагаться на отдельных устройствах, удаленных локациях или даже в~различных облачных средах. Компоненты взаимодействуют друг с другом посредством сетевых вызовов, сообщений или использования распределенных протоколов\cite{ifmo2024}.

Такой подход позволяет повысить надёжность системы, упростить модернизацию отдельных компонентов и обеспечить масштабируемость и~безопасность ИС\cite{irgups2025}.

Архитектура информационной системы является фундаментом её успешного функционирования. Чёткое выделение компонентов и продуманное взаимодействие между ними позволяют создавать эффективные, устойчивые и гибкие информационные системы. Понимание принципов архитектуры ИС важно как для специалистов в области информационных технологий, так~и~для~пользователей, принимающих участие в эксплуатации и развитии систем.
\newpage
\section{Роли серверов, клиентов и~сетевых устройств в~архитектуре системы}
Архитектура «клиент-сервер» -- основа большинства современных
информационных систем. В этой модели взаимодействие между клиентами и
серверами происходит через сети, где каждый элемент выполняет определённую
роль. Серверы предоставляют ресурсы, а клиенты их запрашивают, сетевые
устройства обеспечивают бесперебойную связь между ними.
\subsection{Роль серверов в архитектуре «клиент-сервер»}
Сервер в архитектуре клиент-сервер -- это центральный компонент системы,
который обрабатывает запросы клиентов и предоставляет им необходимые
ресурсы. Серверы могут быть разных типов в зависимости от~функций,
которые они выполняют:
\begin{enumerate}[leftmargin=0em, itemindent=3.75em, nosep, label=\arabic*)]
\item Web-серверы обрабатывают HTTP-запросы и предоставляют доступ к~вебконтенту;
\item Серверы баз данных управляют хранением и доступом к данным, что~позволяет
клиентам работать с большими объёмами информации;
\item Файловые серверы предоставляют доступ к централизованным файловым
хранилищам.
\end{enumerate}

Кроме того, серверы могут быть организованы в многослойные архитектуры, где
каждый слой выполняет отдельные функции, улучшая производительность и
масштабируемость системы. В многоуровневой архитектуре серверы могут
быть распределены по типам:
\begin{itemize}[label=-, leftmargin=0em, itemindent=3.35em, nosep]
\item Web-сервер обрабатывает запросы клиентов;
\item Приложенческий сервер управляет бизнес-логикой;
\item Сервер базы данных обеспечивает хранение и обработку данных.
\end{itemize}

В распределённой архитектуре серверы могут взаимодействовать друг с~другом,
повышая отказоустойчивость и устойчивость системы к нагрузкам.
\subsection{Роль клиентов в архитектуре «клиент-сервер»}
Клиенты -- это устройства или программы, которые обращаются к~серверу для
получения информации или ресурсов. Клиенты могут быть разнообразными:
настольные компьютеры, мобильные устройства или специализированные
программы, которые отправляют запросы на сервер.

Существует несколько типов клиентов:
\begin{itemize}[label=-, leftmargin=0em, itemindent=3.35em, nosep]
\item Тонкие клиенты -- минимально взаимодействуют с сервером, например, веббраузеры;
\item Толстые клиенты -- приложения, которые имеют значительную часть бизнеслогики на своей стороне и минимально зависят от сервера.
\end{itemize}

Примеры таких клиентов включают браузеры для веб-сайтов, мобильные
приложения и настольные программы для работы с корпоративными
информационными системами.
\subsection{Роль сетевых устройств в архитектуре «клиент-сервер»}
Сетевые устройства играют критически важную роль в архитектуре клиент-сервер,
обеспечивая передачу данных между клиентами и серверами. Сетевые
устройства, такие как маршрутизаторы, коммутаторы и точки доступа, управляют
потоком трафика, маршрутизируют данные между сетями и~поддерживают
бесперебойную работу системы.

Маршрутизаторы обеспечивают связь между различными сетями и~корректную
маршрутизацию данных.

Коммутаторы распределяют трафик внутри локальной сети, связывая клиентов и
серверы.

Брандмауэры защищают систему от внешних угроз, контролируя входящий и
исходящий трафик.

Эффективность и безопасность сетевых устройств имеют решающее значение
для обеспечения стабильности работы системы в целом.
\subsection{Типы клиент-серверных архитектур}
Архитектура «клиент-сервер» может быть реализована по-разному в~зависимости
от уровня сложности системы и распределения задач.

Одноуровневая архитектура -- используется в простых системах, где~клиент и
сервер взаимодействуют напрямую.

Двухуровневая архитектура -- сервер разделяется на два компонента: один
управляет веб-запросами, другой -- базой данных.

Многоуровневая архитектура -- задачи разделяются между несколькими
серверами, что позволяет гибко масштабировать систему и повышает
отказоустойчивость.
\subsection{Архитектура компьютерных сетей}
Для стабильного взаимодействия клиентов и серверов используются сетевые
протоколы, например TCP/IP, которые обеспечивают безопасную передачу
данных. Для защиты информации через интернет применяются VPN.
Протоколы REST и SOAP позволяют клиентам и серверам обмениваться данными
независимо от платформы и языка программирования.
\subsection{Архитектура надстраиваемых приложений}
Надстраиваемые приложения упрощают создание распределённых систем, где
данные передаются по общим правилам. Протоколы SOAP и~REST
стандартизируют обмен информацией, что ускоряет подключение новых
технологий.
\subsection{Технология клиент-сервер в информационных системах учебных
заведений}
Архитектура клиент-сервер применяется в образовательных учреждениях
для построения систем управления обучением, электронных журналов, баз
данных студентов и дистанционного обучения. Серверы централизуют
хранение данных, клиенты получают доступ через веб-интерфейсы или специализированные приложения.

Применение клиент-серверной архитектуры повышает эффективность работы
образовательных учреждений и облегчает управление процессами, однако
требует соблюдения мер безопасности для защиты чувствительных данных.
\newpage
\section{Функции пользователей, администраторов и операторов в работе информационной системы}
\subsection{Функции пользователей}
Пользователь информационной системы -- лицо, группа лиц или организация, пользующееся услугами информационной системы для~получения информации или решения других задач. Пользователи являются конечными потребителями функционала информационной системы (ИС), а их права соответствуют принципу минимальных привилегий\cite{skyeng_user}.

Роли пользователей в части доступа к ресурсам идентичны бизнес-ролям (должностям) сотрудников, поэтому их использование для~разграничения полномочий позволяет выстраивать стратегию управления доступом в строгом соответствии с корпоративной иерархией, т.е. с~организационно штатной структурой.

Роли могут быть простыми или составными. Простые разрешают какое-то одно действие, например, только чтение файлов, составные -- совмещают полномочия разного уровня. Одному пользователю может быть назначено несколько ролей, если это необходимо для выполнения бизнес-задач. Такой подход позволяет предоставлять сотрудникам одной должности разные привилегии, благодаря чему реализуется гибкий подход к управлению доступом\cite{studfile_user}.
\subsection{Случайный пользователь}
Определение: лицо, эпизодически взаимодействующее с ИС без наличия специальных знаний и не в силу служебных обязанностей.

Задачи такого пользователя обычно ограничиваются заполнением простых форм, поиском справочной информации или подачей разовых заявок через упрощённый интерфейс. Он не обладает глубокими знаниями о системе и не несёт ответственности за её данные, а его доступ ограничен выполнением конкретного действия\cite{nag_it}.
\subsection{Конечный пользователь (потребитель информации)}
Определение: лицо или коллектив, в интересах которых работает ИС. Он работает с ИС повседневно, связан с ограниченной областью деятельности и,~как правило, не является программистом, например, это может быть бухгалтер, экономист, руководитель подразделения.

Функции сосредоточены на операционном вводе и обработке данных: проведении документов, формировании отчётов, управлении задачами и~анализе информации в рамках своих профессиональных компетенций. Именно конечные пользователи потребляют основную часть информационных ресурсов системы и напрямую влияют на достоверность её данных, строго следуя установленным бизнес-регламентам\cite{finam_user}.
\subsection{Коллектив специалистов (персонал ИС)}
Определение: группа сотрудников, осуществляющих проектирование, разработку, внедрение, сопровождение, развитие ИС, обеспечивающих её корректную работу, целостность и безопасность на~техническом, функциональном уровнях.

В коллектив специалистов входят администраторы баз данных, отвечающие за проектирование структур хранения информации, оптимизацию производительности и организацию резервного копирования. Системные аналитики выполняют критически важную роль связующего звена между бизнес-требованиями и технической реализацией: они формализуют потребности конечных пользователей, проектируют бизнес-процессы и~создают технические спецификации для программистов.

Системные программисты обеспечивают стабильность фундаментальных компонентов ИС -- разрабатывают и поддерживают операционные системы, драйверы и серверные службы, обеспечивая бесперебойную работу платформы в целом. Прикладные программисты, в свою очередь, создают и дорабатывают конкретный функционал для конечных пользователей: реализуют бизнес-логику, разрабатывают отчёты и интерфейсы, руководствуясь техническими заданиями аналитиков\cite{nag_it}.
\subsection{Функции администраторов}
Администраторы информационных сетей -- это привилегированные пользователи, наделенные расширенными полномочиями при работе с~производственной IT-инфраструктурой. Они получают специальный доступ к настройкам, установке компонентов и файловых систем, аудиту, обслуживанию баз данных. Соответственно, у администраторов ИС не~обычные, а привилегированные учетные записи.

Администратор, отвечающий за функционирование, сопровождение и~контроль работы ИС:
\begin{enumerate}[leftmargin=0em, itemindent=3.75em, nosep, label=\arabic*)]
\item Согласовывает вопросы интеграции ИС с поставщиками ПО и~аппаратной части;
\item Решает вопросы совершенствования инфраструктуры компании, регулирует процессы интеграции новых сервисов;
\item Подключает новых пользователей, обучает сотрудников работе с~системами;
\item Проводит диагностику систем, собирает статистические данные;
\item Выявляет и устраняет ошибки в функционировании системных и~аппаратных средств;
\item Принимает меры по противостоянию угрозам несанкционированного доступа;
\item Участвует в подготовке инструктажей, справок, нормативных документов.
\end{enumerate}

На масштабном предприятии может быть несколько администраторов информационных систем, распределенным по службам управления, развития, эксплуатации и контроля, общего управления.

В функции службы управления входит инсталляция операционной системы (ОС) и системы управления базами данных (СУБД), задание параметров запуска приложений. Специалисты службы следят за изменениями в работе ос, СУБД и приложений; осуществляют мониторинг и сбор статистики параметров ИС с помощью средств диагностики ошибок. Осуществляют комплекс мер по противодействию различным угрозам несанкционированного доступа. Службу управления для крупных предприятий рекомендовано разделить на: службу управления конфигурацией, службу контроля характеристик, службу управления ошибочными ситуациями, службу безопасности и службу управления производительностью ИС.

Служба планирования и развития определяет техническую эффективность от внедрения различного вида информационных услуг или сервисов компании, следят за появлением новых технологий и оценивают целесообразность их использования. Контролируют выполнение подрядчиками работ по внедрению этих технологий.

Служба эксплуатации и контроля осуществляет архивирование и~восстановление информационной системы. Ведут базу данных копий, определяют стратегию восстановления ИС в случае ее повреждения. Они же занимаются сопровождением аппаратных средств (например, заменой картриджа принтера), подключением новых пользователей, выполнение профильных работ (уход за оборудованием).

Служба общего управления контролирует работу всех перечисленных служб, следит за согласованием их действий, разрабатывает инструкции для пользователей, внедряет нормативную и справочную документацию, необходимую предприятию\cite{studfile_user}.
\subsection{Функции операторов}
Оператор информационных систем и ресурсов -- это специалист, который обеспечивает функционирование информационно-технической инфраструктуры организации. Он поддерживает работоспособность систем хранения, обработки и передачи данных, следит за их безопасностью и~доступностью для пользователей\cite{sky_operator}. 

В отличие от системных администраторов, которые в основном занимаются настройкой и обслуживанием технических компонентов, операторы ИС фокусируются на информационной составляющей.

Задачи оператора информационных систем стоят на пересечении нескольких областей:
\begin{enumerate}[leftmargin=0em, itemindent=3.75em, nosep, label=\arabic*)]
\item Информационные технологии (управление системами и базами данных).
\item Информационная безопасность (контроль доступа к данным).
\item Бизнес-процессы (понимание потребностей компании в информации).
\item Пользовательская поддержка (помощь сотрудникам в работе с~информационными ресурсами)\cite{sky_operator}.
\end{enumerate}

В зависимости от размера организации и её специфики, обязанности оператора ИС могут варьироваться от базового обслуживания информационных систем до участия в их разработке и оптимизации.

Задачи оператора информационных систем и ресурсов охватывают широкий спектр действий, связанных с обеспечением бесперебойной работы информационной инфраструктуры компании. Основные направления деятельности этих специалистов:
\begin{enumerate}[leftmargin=0em, itemindent=3.75em, nosep, label=\arabic*)]
\item Администрирование информационных систем;
\item Управление данными и информационными ресурсами;
\item Обеспечение информационной безопасности;
\item Техническая поддержка пользователей;
\item Аналитическая работа\cite{sky_operator}.
\end{enumerate}

Под администрированием информационных систем обычно понимают поддержание работоспособности баз данных и прикладного программного обеспечения. Эта деятельность включает мониторинг производительности систем и предупреждение сбоев.

В случае с управлением данными и информационными ресурсами ключевой задачей становится организация хранения и структурирования данных. Основой для этого служит обеспечение целостности и актуальности информации. Обязательным компонентом управления является создание резервных копий и восстановление данных при необходимости.

Обеспечение информационной безопасности состоит из управления правами доступа пользователей к информационным ресурсам. Основу защиты составляет реализация политик безопасности в части информационных систем. 

Техническая поддержка пользователей -- консультирование по вопросам работы с информационными системами. Центральное место занимает устранение проблем доступа к данным и информационным ресурсам. Важным направлением деятельности является создание инструкций и обучающих материалов по работе с ИС. 

Аналитическая работа состоит из формирования отчётов о работе информационных систем. Далее -- анализ использования информационных ресурсов компании. На основе этого анализа происходит выявление узких мест и предложение решений по оптимизации. Перспективное развитие обеспечивается исследованием новых технологий для повышения эффективности работы ИС\cite{sky_operator}.
\newpage
\section{Принципы организации доступа к информационным ресурсам}
В условиях цифровизации и широкого использования информационных 
технологий информация стала одним из ключевых ресурсов любой организации. 
От правильного хранения, обработки и защиты информации напрямую зависит 
эффективность работы предприятий, образовательных учреждений и 
государственных структур. При этом особое значение имеет защита информации 
от несанкционированного доступа, утечки и искажения.

Одним из основных элементов системы информационной безопасности является 
организация доступа к информационным ресурсам. Она определяет правила, 
условия и способы взаимодействия пользователей с информацией. Грамотно 
выстроенная система доступа позволяет ограничить действия пользователей в 
рамках их должностных обязанностей и тем самым повысить уровень 
защищённости данных. 
\subsection{Понятие и роль управления доступом }
Управление доступом к информационным ресурсам представляет собой 
совокупность организационных и технических мер, направленных на~регулирование прав пользователей на получение, изменение и использование 
информации. Доступ может предоставляться как сотрудникам организации, так и 
автоматизированным системам. 

Информационные ресурсы включают в себя электронные документы, базы 
данных, файловые серверы, корпоративные информационные системы и~другие 
виды данных. Без чётко определённых правил доступа возрастает риск утечки 
конфиденциальной информации и нарушения целостности данных. 

Таким образом, управление доступом является неотъемлемой частью системы 
информационной безопасности и обеспечивает защиту информации на уровне 
пользователей. 
\subsection{Принцип идентификации и аутентификации}
Идентификация и аутентификация являются основой любой системы управления 
доступом. Идентификация позволяет системе определить пользователя по 
уникальному признаку, например имени учетной записи. Аутентификация 
подтверждает, что пользователь действительно является тем, за кого себя выдаёт.

Для аутентификации могут использоваться пароли, электронные ключи, смарт
карты и другие средства. Реализация данного принципа позволяет исключить 
доступ посторонних лиц к информационным ресурсам. 
\subsection{Принцип разграничения прав доступа}
Принцип разграничения доступа в разделении прав пользователей в~зависимости от их обязанностей. Пользователю предоставляется доступ только к 
тем информационным ресурсам, которые необходимы для выполнения его 
служебных функций. 

Этот принцип снижает вероятность случайного или умышленного нарушения 
Принцип учета и контроля действий пользователей 
безопасности, так как даже при компрометации учетной записи злоумышленник 
не сможет получить доступ ко всей информации организации. 
\subsection{Принцип минимальных привилегий} 
Принцип минимальных привилегий тесно связан с разграничением доступа и 
предполагает предоставление пользователю необходимого набора 
прав. Это позволяет сократить потенциальный ущерб в случае ошибок 
пользователей или атак. 

Применение данного принципа важно в крупных информационных 
системах с большим количеством пользователей. 
\subsection{Принцип ролевого управления доступом}
Ролевое управление доступом предполагает назначение пользователям 
определённых ролей, каждая из которых имеет заранее определённый набор 
прав. Например, роли могут соответствовать должностям или функциям 
сотрудников. 

Использование ролей упрощает администрирование системы доступа, позволяет 
быстрее изменять права пользователей и повышает прозрачность управления 
доступом. 

Для обеспечения информационной безопасности важно вести учет действий 
пользователей в системе. Принцип учета и контроля предполагает регистрацию 
событий доступа, попыток входа и операций с информацией. 

Анализ журналов позволяет выявлять нарушения, подозрительные действия и 
своевременно реагировать на инциденты информационной безопасности. 
\subsection{Принцип централизованного управления доступом}
Централизованное управление доступом предполагает использование единого 
механизма или системы для контроля прав доступа ко всем информационным 
ресурсам организации. Это снижает вероятность ошибок и~упрощает 
администрирование. 

Централизованный подход позволяет быстро блокировать доступ уволенных 
сотрудников и оперативно вносить изменения в права пользователей. 
\subsection{Практическое значение организации доступа} 
Применение принципов организации доступа к информационным ресурсам 
способствует повышению уровня информационной безопасности организации. 
Четко регламентированный доступ снижает риск утечки информации, связанный 
с человеческим фактором. 

Кроме того, управление доступом помогает соблюдать внутренние нормативные 
документы и требования информационной безопасности, а~также упрощает 
контроль за использованием информационных ресурсов. 

Организация доступа к информационным ресурсам является важным элементом 
системы информационной безопасности. Основные принципы, такие как 
идентификация и аутентификация, разграничение и минимизация прав доступа, 
ролевое управление, учет и контроль действий пользователей, а также 
централизованное управление, обеспечивают защиту информации от~несанкционированного доступа. 

Реализация данных принципов снижает риски информационных угроз и 
обеспечить безопасное использование информационных ресурсов в~организации. 
\newpage
\section{Основные задачи системного администрирования и поддержания работоспособности ИС}
В современном мире информационные системы (ИС) являются кровеносной системой любого предприятия -- от небольшого офиса до~крупной корпорации. Бесперебойная и безопасная работа ИС напрямую влияет на эффективность бизнес-процессов. Ключевую роль в обеспечении этой работоспособности играет системный администратор.
\subsection{Понятие системного администрирования и информационной системы}
Системное администрирование -- это комплекс мероприятий по~контролю, управлению,  обслуживанию и развитию ИС, включая её аппаратную часть, программное обеспечение и сетевую инфраструктуру.

Информационная система (ИС) -- это взаимосвязанная совокупность средств, методов и персонала, используемых для хранения, обработки и~выдачи информации в интересах достижения поставленной цели.
\subsection{Основные задачи системного администрирования}
Задачи системного администратора можно разделить на несколько ключевых блоков.

Установка, настройка и сопровождение ПО и оборудования:
\begin{enumerate}[leftmargin=0em, itemindent=3.75em, nosep, label=\arabic*)]
\item Развёртывание ОС и серверного ПО: установка операционных систем (Windows Server, Linux), настройка служб (Active Directory, DNS, DHCP, веб-серверов, СУБД);
\item Управление оборудованием: инвентаризация, сборка, настройка и~ремонт серверов, рабочих станций, сетевых устройств (маршрутизаторов, коммутаторов);
\item Обновление ПО: планирование и применение патчей, обновлений безопасности и версионных обновлений приложений.
\end{enumerate}

Обеспечение сетевой безопасности и управление доступом:
\begin{itemize}[label=-, leftmargin=0em, itemindent=3.35em, nosep]
\item Аутентификация и авторизация: создание и управление учётными записями пользователей и групп, назначение прав доступа к ресурсам (файлам, папкам, принтерам);
\item Защита периметра: настройка и обслуживание межсетевых экранов (фаерволов), систем обнаружения и предотвращения вторжений (IDS/IPS);
\item Борьба с угрозами: внедрение и поддержка антивирусного ПО, защита от вредоносного ПО, фишинга;
\item Резервное копирование и аварийное восстановление (DR): организация регулярного бэкапа критически важных данных, проверка их целостности, планирование и тестирование процедур восстановления после сбоев.
\end{itemize}

Мониторинг, анализ и устранение неисправностей:
\begin{enumerate}[leftmargin=0em, itemindent=3.75em, nosep, label=\arabic*)]
\item Проактивный мониторинг: постоянный контроль за ключевыми параметрами системы: загрузка процессора и памяти, свободное место на дисках, сетевая активность, состояние служб. Использование систем мониторинга (Zabbix, Nagios, PRTG).
\item Анализ журналов (Log Management): сбор и анализ логов систем и~приложений для выявления аномалий, следов несанкционированного доступа, ошибок.
\item Техническая поддержка пользователей (Help Desk/Service Desk): оперативное реагирование на инциденты и запросы пользователей, диагностика и решение проблем.
\end{enumerate}

Планирование и развитие ИС:
\begin{itemize}[label=-, leftmargin=0em, itemindent=3.35em, nosep]
\item Емкостное планирование: прогнозирование роста нагрузки на систему, планирование модернизации оборудования и увеличения ресурсов (дискового пространства, оперативной памяти);
\item Документирование: ведение технической документации -- схемы сети, пароли, конфигурации оборудования, инструкции по устранению типовых неисправностей;
\item Разработка и внедрение IT-политик: создание регламентов по~работе с~данными, использованию паролей, использованию интернета и т.д.
\end{itemize}
\subsection{Поддержание работоспособности ИС как ключевая цель}
Все перечисленные задачи подчинены одной главной цели -- обеспечению работоспособности ИС, которая характеризуется тремя принципами:
\begin{enumerate}[leftmargin=0em, itemindent=3.75em, nosep, label=\arabic*)]
\item Доступность (Availability): система и данные должны быть доступны авторизованным пользователям в согласованное время (измеряется в процентах uptime, например, 99.9\%).
\item Целостность (Integrity): гарантия точности и непротиворечивости данных, защита от несанкционированного изменения или уничтожения.
\item Конфиденциальность (Confidentiality): обеспечение доступа к~информации только авторизованным лицам.
\end{enumerate}

Для поддержания работоспособности используется цикл PDCA (Plan-Do-Check-Act):
\begin{itemize}[label=-, leftmargin=0em, itemindent=3.35em, nosep]
\item Plan (Планируй): разработка политик и процедур;
\item Do (Выполняй): внедрение и эксплуатация;
\item Check (Проверяй): мониторинг и аудит;
\item Act (Воздействуй): внесение корректировок на основе результатов проверки.
\end{itemize}

Таким образом, системное администрирование -- это многогранная и~динамичная область, выходящая далеко за рамки «починки компьютеров». Системный администратор выступает в роли архитектора, строителя и~службы безопасности информационной инфраструктуры одновременно. От~эффективного выполнения его основных задач -- от установки оборудования до планирования развития и обеспечения безопасности -- напрямую зависит стабильность, безопасность и конкурентоспособность всей организации. Грамотное администрирование превращает ИС из набора технологий в~надёжный инструмент для достижения бизнес-целей.
\newpage
\section{Инструменты и практики администрирования в~современных вычислительных средах}
\subsection{Понятие вычислительной среды}
	Вычислительная среда -- это совокупность аппаратных и программных средств, обеспечивающих выполнение задач обработки данных. В современных условиях она может включать:
    \begin{itemize}[label=-, leftmargin=0em, itemindent=3.35em, nosep]
\item серверы и рабочие станции;
\item сетевое оборудование;
\item операционные системы;
\item виртуальные машины;
\item облачные сервисы;
\item базы данных;
\item прикладное программное обеспечение.
\end{itemize}

Современные вычислительные среды часто являются гибридными, то~есть объединяют локальную инфраструктуру и облачные решения.
\subsection{Роль системного администратора}
	Системный администратор отвечает за:
        \begin{itemize}[label=-, leftmargin=0em, itemindent=3.35em, nosep]
\item установку и настройку программного обеспечения;
\item поддержку серверов и сетей;
\item управление пользователями и правами доступа;
\item резервное копирование данных;
\item мониторинг состояния систем;
\item обеспечение информационной безопасности.
\end{itemize}

	Сегодня администратор -- это не просто «человек, который чинит компьютеры», а специалист, управляющий сложной цифровой экосистемой.
    
	Пример: в компании с 200 сотрудниками администратор следит за~работой серверов, поддерживает корпоративную почту, настраивает VPN для удалённых работников и контролирует безопасность сети.
\subsection{Операционные системы и средства управления}
Операционные системы являются основой любой вычислительной среды. Наиболее распространены:
\begin{enumerate}[leftmargin=0em, itemindent=3.75em, nosep, label=\arabic*)]
\item Windows Server -- часто используется в корпоративных сетях;
\item Linux (Ubuntu Server, CentOS, Debian) -- популярен благодаря стабильности и гибкости.
\end{enumerate}

	Для администрирования используются встроенные инструменты:
    \begin{itemize}[label=-, leftmargin=0em, itemindent=3.35em, nosep]
\item командная строка (PowerShell, Bash);
\item графические панели управления;
\item журналы событий.
\end{itemize}

Пример: администратор с помощью PowerShell может массово создать учётные записи сотрудников и автоматически выдать им права доступа.
\subsection{Средства удалённого администрирования}
	Удалённый доступ позволяет управлять системами без физического присутствия. Это особенно важно при распределённых и облачных средах.
    
	Популярные инструменты:
    \begin{itemize}[label=-, leftmargin=0em, itemindent=3.35em, nosep]
\item Remote Desktop;
\item SSH;
\item TeamViewer;
\item AnyDesk;
\item VNC.
\end{itemize}

	Пример: администратор подключается по SSH к серверу в дата-центре и~перезапускает службу базы данных после сбоя, не выходя из офиса.
\subsection{Виртуализация и контейнеризация}
	Виртуализация позволяет запускать несколько виртуальных серверов на~одном физическом оборудовании.
	Основные инструменты:
    \begin{enumerate}[leftmargin=0em, itemindent=3.75em, nosep, label=\arabic*)]
\item VMware;
\item Hyper-V;
\item VirtualBox;
\item KVM.
\end{enumerate}

	Контейнеризация (Docker, Kubernetes) облегчает развертывание и~масштабирование приложений.
    
	Пример: на одном сервере работают виртуальные машины для~сайта компании, бухгалтерии и системы видеонаблюдения -- каждая изолирована и~не мешает другим.
\subsection{Облачные платформы}
	Облачные вычисления стали неотъемлемой частью современных сред.
    
	Основные платформы:
    \begin{enumerate}[leftmargin=0em, itemindent=3.75em, nosep, label=\arabic*)]
\item Amazon Web Services;
\item Microsoft Azure;
\item Google Cloud.
\end{enumerate}

	Инструменты облачного администрирования позволяют:
    \begin{itemize}[label=-, leftmargin=0em, itemindent=3.35em, nosep]
\item быстро создавать серверы;
\item управлять хранилищами;
\item масштабировать ресурсы;
\item контролировать затраты.
\end{itemize}

	Пример: интернет-магазин увеличивает мощности сервера в период распродаж, а после снижает их, экономя деньги.
\subsection{Управление пользователями и доступом}
Контроль доступа -- ключевая задача администратора. Используются:
\begin{enumerate}[leftmargin=0em, itemindent=3.75em, nosep, label=\arabic*)]
\item доменные службы;
\item роли пользователей;
\item многофакторная аутентификация.
\end{enumerate}

	Пример: обычный сотрудник может просматривать документы, а~бухгалтер -- дополнительно редактировать финансовые отчёты.
\subsection{Резервное копирование и восстановление}
Резервное копирование защищает данные от потерь из-за сбоев, ошибок или атак. Виды резервного копирования:
\begin{itemize}[label=-, leftmargin=0em, itemindent=3.35em, nosep]
\item полное;
\item инкрементное (сохраняет изменения после любого последнего копирования);
\item дифференциальное (сохраняет изменения после последнего полного копирования).
\end{itemize}

Пример: каждую ночь система автоматически сохраняет копию базы данных на удалённый сервер. В случае сбоя данные можно восстановить за~несколько минут.
\subsection{Мониторинг и журналирование}
	Мониторинг позволяет заранее выявлять проблемы. Используемые инструменты:
    \begin{itemize}[label=-, leftmargin=0em, itemindent=3.35em, nosep]
\item Zabbix;
\item Nagios;
\item Prometheus;
\item Grafana.
\end{itemize}

	Журналы событий помогают анализировать причины сбоев.
	
    Пример: система мониторинга сообщает администратору о перегрузке сервера ещё до того, как пользователи заметят замедление работы.
\subsection{Обновление и патч-менеджмент}
	Регулярные обновления устраняют уязвимости и повышают стабильность. Практики включают:
    \begin{enumerate}[leftmargin=0em, itemindent=3.75em, nosep, label=\arabic*)]
\item тестирование обновлений;
\item поэтапное внедрение;
\item автоматизацию установки патчей.
\end{enumerate}

	Пример: перед обновлением сервера администратор проверяет новую версию на тестовой машине, чтобы избежать сбоев.
\subsection{Информационная безопасность как часть администрирования}
	Современные вычислительные среды подвержены:
    \begin{itemize}[label=-, leftmargin=0em, itemindent=3.35em, nosep]
\item вирусам;
\item фишинговым атакам;
\item утечкам данных;
\item внутренним угрозам.
\end{itemize}

	Для защиты используются:
    \begin{enumerate}[leftmargin=0em, itemindent=3.75em, nosep, label=\arabic*)]
\item антивирусы;
\item межсетевые экраны;
\item системы обнаружения вторжений;
\item VPN.
\end{enumerate}

	Пример: сотрудники подключаются к корпоративной сети из дома только через VPN, что защищает данные от перехвата.
\subsection{Автоматизация администрирования}
	Автоматизация снижает нагрузку на администратора и уменьшает число ошибок. Инструменты:
    \begin{itemize}[label=-, leftmargin=0em, itemindent=3.35em, nosep]
\item PowerShell;
\item Bash;
\item Ansible;
\item Puppet;
\item Chef.
\end{itemize}

	Пример: скрипт автоматически устанавливает программное обеспечение на 50 компьютеров за несколько минут.
	
    DevOps объединяет разработку и администрирование, ускоряя внедрение новых решений. Ключевые идеи:
    \begin{enumerate}[leftmargin=0em, itemindent=3.75em, nosep, label=\arabic*)]
\item автоматизация;
\item непрерывная интеграция;
\item мониторинг в реальном времени.
\end{enumerate}

	Пример: обновление веб-приложения происходит без остановки сервиса для пользователей.
\newpage
{
\sloppy
\nocite{*}
\cleardoublepage
\addcontentsline{toc}{section}{Список использованных источников}
\bibliographystyle{plainurl}
\bibliography{report11}
}
\end{document}